\subsection{Nagata rings: completion, finite fields and the original monogenic rings}\label{algebra-proof-067-nagata-rings-completion-finite-fields-and-the-original-monogenic-rings}

Source: Stacks Project authors, algebra.\AlgebraIdentifierBreak{}tex:\AlgebraIdentifierBreak{}45615--\AlgebraIdentifierBreak{}46292, authority a044\allowbreak{}46e5\allowbreak{}7ec1\allowbreak{}fbc2\allowbreak{}52a8\allowbreak{}71af\allowbreak{}cec7\allowbreak{}752f\allowbreak{}b280\allowbreak{}7b14, SHA256 FA8B\allowbreak{}B92E\allowbreak{}58A4\allowbreak{}F78A\allowbreak{}2BD0\allowbreak{}1B3B\allowbreak{}6A4A\allowbreak{}87DE\allowbreak{}0A0D\allowbreak{}279F\allowbreak{}5DD9\allowbreak{}0641\allowbreak{}B574\allowbreak{}DD5F\allowbreak{}BFFF\allowbreak{}A4F3. The twenty-two received reports are checked against the original source and the ten existing units MC-\AlgebraIdentifierBreak{}STK-\AlgebraIdentifierBreak{}ERR-\AlgebraIdentifierBreak{}0715--\AlgebraIdentifierBreak{}0724. The following complete receiving proofs are editorial material. The source's structure, rings, generators and arguments remain identifiable. Optional extensions and further paper mining remain after Algebra.

\subsubsection{Reduced algebras and nonzero localization charts}\label{algebra-proof-067-reduced-algebras-and-nonzero-localization-charts}

The source defines a universally Japanese ring by requiring every finite-type algebra which is a domain to be N-2. A Nagata ring is Noetherian and has N-2 quotients by all prime ideals. The zero ring satisfies both definitions: it is Noetherian, has no prime ideals, and admits no unital map to a nonzero domain. These observations also cover empty spectra and empty finite covers below.

For 45639--45667 let \(R\) be Nagata and \(S\) an essentially finite-type reduced \(R\)-algebra. It is Noetherian. If its minimal primes are \(\mathfrak q_1,\ldots,\mathfrak q_m\), the maps to its minimal-prime localizations give an injection \[
 S\longrightarrow\prod_{i=1}^m K_i,
 \qquad K_i=S_{\mathfrak q_i}=\operatorname{Frac}(S/\mathfrak q_i).
\] Its kernel is the intersection of the minimal primes, which is zero; each localization is a field because it is reduced, local and zero-dimensional. Put \(\mathfrak p_i=\mathfrak q_i\cap R\). Each \(K_i\) is a finitely generated field extension of \(\kappa(\mathfrak p_i)=\operatorname{Frac}(R/\mathfrak p_i)\). The exact algebraic subfield \(L_i\subset K_i\) is finite over that base field by fields.\AlgebraIdentifierBreak{}tex:\AlgebraIdentifierBreak{}3583--\AlgebraIdentifierBreak{}3603. To recall its degree argument, choose an actual transcendence basis \(t_1,\ldots,t_r\) with \([K_i:\kappa(\mathfrak p_i)(t_1,\ldots,t_r)]=n\). For every finite algebraic subextension \(k'/\kappa(\mathfrak p_i)\) inside \(K_i\), adjoining those same independent variables preserves the original finite degree, so \[
 [k':\kappa(\mathfrak p_i)]
 =[k'(t_1,\ldots,t_r):\kappa(\mathfrak p_i)(t_1,\ldots,t_r)]\leq n.
\] The degree preservation uses the same basis of \(k'\) after rational extension: clear rational denominators in a putative relation, and compare the coefficients of the independent monomials. Algebraic extension of the constant field preserves independence of the variables. Among finite subextensions choose one of maximal degree; adjoining any further algebraic element cannot increase that degree, so it already contains the whole algebraic subfield. This proves finiteness.

An element of \(K_i\) integral over \(R\) is integral over its image \(R/\mathfrak p_i\), and therefore algebraic over \(\kappa(\mathfrak p_i)\). Thus its closure is precisely \(C_i=\operatorname{cl}_{R/\mathfrak p_i}(L_i)\), finite over \(R/\mathfrak p_i\) by the Nagata definition. The closure \(C\) of \(R\) in the original \(S\) maps injectively into \(\prod_i C_i\). Since \(R\) is Noetherian, this submodule of a finite module is finite. If \(S=0\), the product is empty and the closure is zero, also finite.

For the N-1 test at 45669--45688, let \(S\) be an original finite-type domain over \(R\), and let \(L/\operatorname{Frac}(S)\) be finite. Choose finitely many field generators \(\alpha_j\) of \(L\). For their monic polynomials \(P_j(T)=T^{d_j}+\sum_{i=1}^{d_j}a_{j,i}T^{d_j-i}\), choose actual nonzero \(c_j\in S\) clearing all coefficient denominators. Retain \(\beta_j=c_j\alpha_j\) and \[
 c_j^{d_j}P_j(T/c_j)
 =T^{d_j}+\sum_{i=1}^{d_j}c_j^i a_{j,i}T^{d_j-i},
 \qquad \alpha_j=c_j^{-1}\beta_j.
\] Then \(S'=S[\beta_1,\ldots,\beta_r]\subset L\) is finite over \(S\), has fraction field exactly \(L\), and is finite type over \(R\). Under the assumed N-1 test its closure \(S''\) in \(L\) is finite over \(S'\), hence over \(S\). Integrality in both directions identifies \(S''=\operatorname{cl}_S(L)\). This proves the N-2 property for the original \(S\), without a Noetherian hypothesis on \(R\).

The essentially finite-type assertion at 45690--45700 has the following exact descent. Write the original algebra as \(A=W^{-1}B\) with \(B\) finite type over a universally Japanese \(R\). Given an arbitrary finite-type domain \(T\) over \(A\), let \(C\subset T\) be the ring generated by the image of \(B\) and finitely many chosen \(A\)-algebra generators of \(T\). Then \(C\) is a finite-type domain over \(R\), and the actual fraction map identifies \(T=W^{-1}C\). The images of \(W\) are nonzero because they are units in the domain \(T\). The corrected nonzero-localization theorem from the Japanese-rings note gives N-2 for \(T\). Thus \(A\) is universally Japanese.

For quasi-finite ascent at 45702--45718, \(S\) is Noetherian because it is finite type over the Noetherian \(R\). For \(\mathfrak q\in\operatorname{Spec}(S)\), the exact injection \(R/\mathfrak p\to S/\mathfrak q\), \(\mathfrak p=\mathfrak q\cap R\), is quasi-finite: it is the quotient of the base change to \(R/\mathfrak p\) and has finite fibres. The original domain ascent theorem from the Japanese-rings note applies, since \(R/\mathfrak p\) is N-2. This proves the Nagata assertion. Localization at \(W\) preserves Noetherianity. Every prime in \(W^{-1}R\) comes from a prime \(\mathfrak p\) disjoint from \(W\), and its domain quotient is the nonzero localization \(W^{-1}(R/\mathfrak p)\); it is N-2. If zero belongs to \(W\), the result is the zero Nagata ring, not a domain to which N-2 is assigned.

For 45729--45752 retain \(\sum_i r_i f_i=1\) in \(R\). Given \(\varphi:R\to S\) with \(S\) a finite-type domain, the elements with \(\varphi(f_i)=0\) define zero charts. Remove exactly those indices from the identity \[
 1=\sum_{\varphi(f_i)\ne0}\varphi(r_i)\varphi(f_i).
\] For each surviving index, \(S_{\varphi(f_i)}\) is a finite-type domain over the universally Japanese \(R_{f_i}\), hence N-2 and N-1. The finite-cover argument proves N-1 for \(S\), and the preceding N-1 test proves the first conclusion. Accordingly the new minimal qualification at 45742 is ``each nonzero'' localization. This propagates ALGEBRA-\AlgebraIdentifierBreak{}RECON-\AlgebraIdentifierBreak{}1292 to the first covering assertion.

For the second assertion, Noetherianity descends along this finite standard cover: for every ideal \(I\subset R\), choose a finite set of elements of \(I\) generating each localization \(I_{f_i}\), and let \(I_0\) be their joint span. Every element of \(I/I_0\) is killed by powers of all \(f_i\); those powers still generate the unit ideal, so \(I=I_0\). For a prime \(\mathfrak p\subset R\), keep only indices \(f_i\notin\mathfrak p\). Their images still generate 1 in \(R/\mathfrak p\), and \[
 R_{f_i}/\mathfrak pR_{f_i}=(R/\mathfrak p)_{f_i}
\] is a domain and N-2. The domain cover lemma applies. This is precisely the already composed MC-\AlgebraIdentifierBreak{}STK-\AlgebraIdentifierBreak{}ERR-\AlgebraIdentifierBreak{}0715. Both parts now use nonzero domain charts, while the source's Nagata assertion itself retains the zero ring.

\subsubsection{Complete rings and integral closure in the total quotient ring}\label{algebra-proof-067-complete-rings-and-integral-closure-in-the-total-quotient-ring}

For 45754--45778 let \(R\) be a complete Noetherian local ring. A prime ideal is proper, so its quotient is again a complete Noetherian local domain by the proper-quotient result in the Cohen note. The original finite regular-subring theorem embeds either \(k[[X_1,\ldots,X_d]]\) or \(\Lambda[[X_1,\ldots,X_d]]\) in that domain, with finite extension. Its two cases, coefficient maps, actual prime \(p\) and all parameters remain those proved in the Cohen note. A field is N-2. For a Cohen ring \(\Lambda\), the actual element \(p\) generates its maximal ideal, its residue quotient is a field, and it is a complete Noetherian normal domain. Tate's theorem from the Japanese-rings note therefore proves N-2 for \(\Lambda\). Iterating that note's power-series theorem preserves all variables and gives N-2 for either regular subring. Finite domain ascent then gives N-2 for the original quotient. Hence \(R\) is Nagata. The empty list of parameters is included.

Write \(B=\widehat R=\varprojlim_n R/\mathfrak m^n\). The source completion theorems give its actual faithfully flat local map, Noetherianity and completeness. A ring is analytically unramified exactly when this \(B\) is reduced. For 45807--45887, faithful flatness makes \(R\to B\) injective, so reducedness of \(B\) implies reducedness of \(R\).

If \(\mathfrak q\) is a minimal prime of such an \(R\), it is associated, so choose the actual \(f\in R\) with annihilator \(\mathfrak q\). Tensoring the kernel sequence of multiplication by \(f\) gives \[
 \operatorname{Ann}_B(f)=\mathfrak qB,
 \qquad \widehat{R/\mathfrak q}=B/\mathfrak qB.
\] If \(z^2\in\mathfrak qB\), then \(fz^2=0\) and \((fz)^2=f(fz^2)=0\). Reducedness of \(B\) implies \(fz=0\), hence \(z\in\mathfrak qB\). This square-root criterion makes the ideal radical: repeated squaring reaches a power at least any given nilpotence exponent and then descends. Thus \(B/\mathfrak qB\) is reduced. Conversely, the actual injection \(R\to\prod_i R/\mathfrak q_i\) for a reduced Noetherian ring remains injective under completion, giving \[
 B\longrightarrow\prod_i\widehat{R/\mathfrak q_i}.
\] If every factor is reduced, so is the subring \(B\).

Assume now that \(B\) is reduced. Its total quotient ring is the product of the fields \(B_{\mathfrak p_i}\) at its finitely many minimal primes. The integral closure \(D\) of \(B\) in that product is \[
 D=\prod_i D_i,\qquad
 D_i=\operatorname{cl}_{B/\mathfrak p_i}\bigl(\operatorname{Frac}(B/\mathfrak p_i)\bigr).
\] For this equality, an integral tuple has integral coordinates. Conversely, lift to \(B[T]\) a monic equation for each coordinate; the product of the finitely many lifted monic polynomials vanishes on the whole tuple, coordinate by coordinate. Since \(B\) is Nagata, each \(D_i\) is finite over \(B/\mathfrak p_i\), hence \(D\) is finite over \(B\).

Keep the original total quotient ring \(Q(R)=U^{-1}R\), with \(U\) the original set of nonzerodivisors, and let \(R'=\operatorname{cl}_R(Q(R))\). For each \(u\in U\), flatness of \(B\) preserves the injection given by multiplication by \(u\). Thus its image is a nonzerodivisor in \(B\). The exact maps are \[
 R'\otimes_R B\ \hookrightarrow\ Q(R)\otimes_R B
 \xrightarrow{\ \cong\ }U^{-1}B\ \hookrightarrow\ Q(B),
 \qquad (r/u)\otimes b\longmapsto rb/u.
\] The first injection is flatness; the last is localization at a subset of the nonzerodivisors. Their image is integral over \(B\), since integrality is preserved by this base change. Hence \(R'\otimes_R B\subset D\), a finite \(B\)-module; Noetherianity makes the submodule finite. Choose module generators and express each as a finite sum \(\sum_j r'_j\otimes b_j\). The finite list of all \(r'_j\) generates the tensor product over \(B\). If \(N\subset R'\) is their \(R\)-span, then \((R'/N)\otimes_R B=0\); faithful flatness implies \(R'=N\). This proves finiteness of the original integral closure, and its domain case is N-1. MC-\AlgebraIdentifierBreak{}STK-\AlgebraIdentifierBreak{}ERR-\AlgebraIdentifierBreak{}0716 corrects the article at 45877; the new copyedit at 45872 supplies ``Denote by''. Neither changes these maps.

\subsubsection{Codimension one and the completion criterion}\label{algebra-proof-067-codimension-one-and-the-completion-criterion}

For 45889--45914, \(R\) need not itself be a domain. Let \(R_{\mathfrak p}\) be the original DVR and suppose \(B/\mathfrak pB\) is reduced. Its associated primes are minimal primes. For one such \(\mathfrak q\), flatness of \(R/\mathfrak p\to B/\mathfrak pB\) and going down show \(\mathfrak q\cap R=\mathfrak p\): if the contraction were larger, a prime below it over zero would contradict minimality in the quotient. Therefore the map \[
 R_{\mathfrak p}\longrightarrow B_{\mathfrak q}
\] is defined and flat. The quotient \((B/\mathfrak pB)_{\mathfrak q}\) is a reduced zero-dimensional local ring, hence a field; thus the maximal ideal of \(B_{\mathfrak q}\) is \(\mathfrak pB_{\mathfrak q}\). Choose \(x\in R\) whose image generates \(\mathfrak pR_{\mathfrak p}\); if an original uniformizer is given as a fraction, its nonzero numerator differs from it by the displayed denominator unit. Flatness of this localized map makes multiplication by \(x\) injective on \(B_{\mathfrak q}\). Its maximal ideal is \(xB_{\mathfrak q}\). A nonzero Noetherian local ring with maximal ideal generated by a nonzerodivisor is a DVR: separatedness of the maximal-ideal powers writes every nonzero element as \(x^n u\) with \(u\) a unit, and the full product \(x^{n+m}uv\) cannot vanish. This gives a domain and the discrete valuation by the actual exponent. No assertion that \(x\) is globally a nonzerodivisor of \(R\) is required.

For the criterion at 45916--45962 let \(R\) be the original Noetherian local domain, \(0\ne x\in\mathfrak m\), and assume the source's three conditions. The associated primes \(\mathfrak p_i\) of \(R/xR\) are its minimal support primes, all of height one by the principal ideal theorem. The regular local rings \(R_{\mathfrak p_i}\) are DVRs. For every \(\mathfrak q_{ij}\in\operatorname{Ass}_B(B/\mathfrak p_iB)\), the preceding paragraph proves that \(B_{\mathfrak q_{ij}}\) is a DVR. Apply the exact flat associated-prime formula of algebra.\AlgebraIdentifierBreak{}tex:\AlgebraIdentifierBreak{}15773--\AlgebraIdentifierBreak{}15851 to \(M=R/xR\) and \(N=B\): \[
 \operatorname{Ass}_B(B/xB)
 =\bigcup_i\operatorname{Ass}_B(B/\mathfrak p_iB)
 =\{\mathfrak q_{ij}\}_{i,j}.
\] Here all hypotheses are present: \(R\) is Noetherian and \(B\) is flat over \(R\). The formula's reverse inclusion follows by tensoring each injection \(R/\mathfrak p_i\to M\). For the forward inclusion, localize an associated prime of the tensor product and contract it to \(R\). If that contraction is not associated to \(M\), a nonzerodivisor in the contracted maximal ideal remains one after the flat tensor map, contradicting association. A finite prime filtration of \(M\), tensored by the flat module, then identifies the corresponding quotient \(B/\mathfrak p_iB\). This is the original proof with these specific inputs.

If \(y\in B\) satisfies \(y^2=0\), it vanishes in every domain \(B_{\mathfrak q_{ij}}\). Consequently its class vanishes in \(B/xB\): a nonzero cyclic submodule of this finite module has an associated prime of the module and cannot vanish at that localization. Thus \(y=xy_1\). Since multiplication by \(x\) on \(B\) is injective by flatness, \(x^2y_1^2=0\) implies \(y_1^2=0\). Iterating gives the actual equalities \(y=x^n y_n\) for all \(n\). The original local Noetherian \(B\) has \(x\) in its maximal ideal, so \(\bigcap_n x^nB=0\) and \(y=0\). A nonzero nilpotent would have a nonzero square-zero power, so \(B\) is reduced.

\subsubsection{Local Nagata domains and the finite local test}\label{algebra-proof-067-local-nagata-domains-and-the-finite-local-test}

For 45964--46015 induct on the finite dimension \(d\) of the original Noetherian local domain \(R\). In dimension zero it is a field. Let \(S\) be its closure in its original fraction field. N-2 for \(R\) makes \(S\) finite; finite ascent makes it Nagata, and it is normal. Its finitely many maximal ideals \(\mathfrak m_i\) all lie over \(\mathfrak m\). Put \(J=\bigcap_i\mathfrak m_i\). Since \(S/\mathfrak mS\) is Artinian, choose the actual integer \(c\geq1\) with \[
 J^c\subseteq\mathfrak mS\subseteq J,
 \qquad J^{cn}\subseteq\mathfrak m^nS\subseteq J^n.
\] The identity induces mutually inverse comparisons of the two completion limits. Chinese remainder and localization of each quotient give \[
 \widehat R\hookrightarrow\widehat S^{\,\mathfrak mS}
 \cong\prod_i\varprojlim_n S/\mathfrak m_i^n
 \cong\prod_i\widehat{S_{\mathfrak m_i}}.
\] The injection is exact completion of the inclusion of finite \(R\)-modules. Every map in the product comparison comes from the same residue or fraction class. The local domains \(S_{\mathfrak m_i}\) are normal Nagata rings of dimension at most \(d\). It suffices to prove reducedness of their completions.

For any such normal local Nagata domain \(A\) of positive dimension at most \(d\), choose the original nonzero \(x\) in its maximal ideal. The height-one intersection result, with its complete proof in the Serre note, gives no embedded primes for \(A/xA\); all associated primes \(\mathfrak p\) have height one and \(A_{\mathfrak p}\) is a DVR. The quotient \(A/\mathfrak p\) is a local Nagata domain of dimension strictly less than \(d\): append \((0)\subsetneq\mathfrak p\) to a chain of primes above \(\mathfrak p\) to obtain the strict bound. No catenary dimension formula is used. Induction makes \(\mathfrak p\) analytically unramified. The preceding criterion now proves that \(A\) is analytically unramified. This treats all the local factors, and their product contains the original completion injectively, proving the original assertion for \(R\).

For 46017--46060, condition (1), that \(R\) is Nagata, implies condition (2) by finite Nagata ascent, localization and the local-domain result just proved. Condition (2) implies (3) by taking the already local finite domain. The missing colon at 46020 and redundant ``then'' at 46026 are copyedits.

Assume (3), and retain an arbitrary original prime \(\mathfrak p\subset R\), \(A=R/\mathfrak p\), and finite \(L/\kappa(\mathfrak p)\). If \(\mathfrak p=\mathfrak m\), \(A\) is a field and its closure in \(L\) is \(L\), a finite module. Otherwise choose the actual \(g\in\mathfrak m\setminus\mathfrak p\) required by MC-\AlgebraIdentifierBreak{}STK-\AlgebraIdentifierBreak{}ERR-\AlgebraIdentifierBreak{}0717. Retain each original field generator \(x_i\) and its minimal polynomial \[
 P_i(T)=T^{d_i}+\sum_{j=1}^{d_i}a_{i,j}T^{d_i-j}.
\] Choose \(f_i\in R\setminus\mathfrak p\) such that all \(\overline f_i a_{i,j}\in A\). With bars denoting actual quotient images, put \[
 y_i=\overline g\,\overline f_i x_i,\qquad
 Q_i(T)=(\overline g\,\overline f_i)^{d_i}
 P_i\bigl(T/(\overline g\,\overline f_i)\bigr)
 =T^{d_i}+\sum_{j=1}^{d_i}
   \overline g^{,j}\overline f_i^{,j}a_{i,j}T^{d_i-j}.
\] All nonleading coefficients belong to \(\mathfrak m/\mathfrak p\). Every multiplier is nonzero in \(A\), and the inverse equation is \(x_i=(\overline g\,\overline f_i)^{-1}y_i\). Thus \(S=A[y_1,\ldots,y_n]\subset L\) is finite over \(R\), a domain, and has the original fraction field \(L\). The special fibre is the image of \[
 (R/\mathfrak m)[T_1,\ldots,T_n]/(T_1^{d_1},\ldots,T_n^{d_n})
 \longrightarrow S/\mathfrak mS,
 \qquad T_i\longmapsto[y_i].
\] This map is onto. The fibre is nonzero by lying over for the integral inclusion \(A\subset S\); the displayed source is local with nilpotent maximal ideal, so its nonzero quotient has one maximal ideal. Every maximal ideal of the finite \(R\)-algebra \(S\) lies over \(\mathfrak m\). Hence \(S\) is local. Condition (3) makes it analytically unramified, so its closure \(S'\) in \(L\) is finite. Transitivity identifies \(S'=\operatorname{cl}_A(L)\), and it is finite over \(A\). This proves Nagata for \(R\).

The new parentheses at 46052 and 46054 identify the original coefficient quotients; they do not change their rings. The suggestion in OCC-01009 to repeat ``in \(L\)'' is an optional clarification, not a mathematical defect: the immediately preceding sentence fixes \(L/\kappa(\mathfrak p)\), the generator construction takes place in that sole extension, and the last sentence explicitly identifies both closures in \(L\). This editorial proof states every ambient field, while the source wording remains. The \(\mathfrak p=\mathfrak m\) case is different: without MC-\AlgebraIdentifierBreak{}STK-\AlgebraIdentifierBreak{}ERR-\AlgebraIdentifierBreak{}0717 its prescribed multiplication could annihilate all generators and lose the given extension.

\subsubsection{Finite-type stability with every original ring retained}\label{algebra-proof-067-finite-type-stability-with-every-original-ring-retained}

For the proposition at 46066--46222, condition (3) implies condition (2) by essentially finite-type stability of universal Japaneseness and Noetherianity of finite-type algebras. Condition (2) implies (1) by the identity algebra. It remains to prove that Nagata implies finite-type stability. Choose actual algebra generators of the original \(R\)-algebra \(S\). Adjoining them successively factors \(R\to S\) into monogenic maps. At each step, and for each prime quotient of the monogenic target, pass to the corresponding prime quotient of the source. Nagata is preserved under that quotient directly from the definition. Thus it suffices to prove N-2 for an original injection of domains \(R\subset S=R[x]\), with \(R\) Nagata.

Write \(K=\operatorname{Frac}(R)\). If \(x\) is transcendental over \(K\), the evaluation \(R[X]\to S\), \(X\mapsto x\), is injective and onto. The polynomial N-2 result applies. Otherwise \(K(x)/K\) is finite. The source's wording ``purely transcendental'' at 46101 needs the qualification ``of positive transcendence degree'' for this first inference: fields.\AlgebraIdentifierBreak{}tex:\AlgebraIdentifierBreak{}3371--\AlgebraIdentifierBreak{}3384 permits a polynomial ring with no variables, so \(K/K\) is itself purely transcendental of degree zero. For example, \(\mathbf Z\subset\mathbf Z[1/2]\) has identical fraction fields, but the evaluation \(\mathbf Z[X]\to\mathbf Z[1/2]\), \(X\mapsto1/2\), has the nonzero relation \(2X-1\). The finite case below handles this example. The qualification selects exactly the transcendental-generator case and leaves all finite cases to the next step.

For the finite case retain \(L=\operatorname{Frac}(S)\), an arbitrary finite extension \(M/L\), and the same original generator \(x\). Define \[
 A=\operatorname{cl}_R(M),\qquad B=A[x]\subset M,
 \qquad C=\operatorname{cl}_S(M).
\] Since \(M/K\) is finite and \(R\) is N-2, \(A\) is finite over \(R\); its fraction field is exactly \(M\) by the multiplier calculation in the Japanese note. It is normal and Nagata by finite ascent. The actual inclusion \(S\subset B\) is finite: a finite \(R\)-generating list of \(A\) also generates \(B\) over \(S\). It is therefore integral, and \[
 C=\operatorname{cl}_B(M).
\] Indeed integrality over \(S\) implies integrality over \(B\), while transitivity gives the reverse implication. Finiteness of \(C\) over \(B\) implies finiteness over \(S\) by the displayed finite inclusion; conversely any finite \(S\)-generating list also generates it over \(B\). Thus it is enough to prove the following unchanged original subproblem: for a normal Nagata domain \(A\), an element \(x\in\operatorname{Frac}(A)\), and \(B=A[x]\) as a subring of that actual fraction field, \(B\) is N-1. No normality of the original \(R\) was needed. This verifies precisely the deletion in MC-\AlgebraIdentifierBreak{}STK-\AlgebraIdentifierBreak{}ERR-\AlgebraIdentifierBreak{}0718, rather than silently restricting the theorem's input.

We now prove that subproblem, initially writing its original normal Nagata domain as \(R\), its actual fraction field as \(K\), and its original generator as \(x=b/a_0\), \(a_0\ne0\). Then \(S=R[x]\subset K\), and the exact fraction maps give \(S_{a_0}=R_{a_0}\), which is normal. The N-1 criterion reduces to each original maximal ideal \(\mathfrak m\subset S\). Let \(\mathfrak n=\mathfrak m\cap R\). The residue field \(\kappa(\mathfrak m)\) is generated as an algebra over \(\kappa(\mathfrak n)\) by the image of \(x\): it is already generated over the subring \(R/\mathfrak n\), and contains its fraction field. By the field version of the Nullstellensatz, this extension is finite. Let \[
 \overline f(X)=X^d+\sum_{i=1}^{d}\overline a_iX^{d-i}
 \in\kappa(\mathfrak n)[X]
\] be its monic minimal polynomial, with the actual coefficients retained. Choose representatives of each coefficient as fractions modulo \(\mathfrak n\), and let \(c\in R\setminus\mathfrak n\) be the product of their chosen nonzero denominator representatives. Each coefficient lifts to \(a_i\in R_c\). Put \[
 f(X)=X^d+\sum_{i=1}^{d}a_iX^{d-i}\in R_c[X].
\] Its evaluation at the unchanged \(x\) belongs to \(\mathfrak mS_c\), because its residue is \(\overline f(\overline x)=0\). Localization preserves the original local ring: \(S_{\mathfrak m}\to(S_c)_{\mathfrak mS_c}\) sends every fraction to the same fraction, with inverse using that \(c\notin\mathfrak m\). The source uses its symbol \(a\) again for this localization element; here \(a_0\) and \(c\) are kept distinct so its two actual denominators cannot be confused. MC-\AlgebraIdentifierBreak{}STK-\AlgebraIdentifierBreak{}ERR-\AlgebraIdentifierBreak{}0719 supplies the missing subject, and MC-\AlgebraIdentifierBreak{}STK-\AlgebraIdentifierBreak{}ERR-\AlgebraIdentifierBreak{}0720 keeps the polynomial indeterminate \(X\) distinct from the element \(x\).

Choose a finite splitting field \(K''/K\) of \(f\), retain all its roots with their multiplicities as \(\rho_1,\ldots,\rho_d\), and set \[
 A=\operatorname{cl}_{R_c}(K''),\qquad B=A[x].
\] Each \(\rho_i\) lies in \(A\) because \(f\) is monic. The ring \(A\) is finite, normal and Nagata over \(R_c\). The original \(S_c\subset B\) is finite. Set \(W=S_c\setminus\mathfrak mS_c\). The ring \(D=W^{-1}B\) is finite over the local domain \((S_c)_{\mathfrak mS_c}=S_{\mathfrak m}\), and its maximal ideals correspond exactly to maximal ideals \(\mathfrak m'\) of \(B\) lying over \(\mathfrak mS_c\). Their local rings are the original \(B_{\mathfrak m'}\). Moreover, \(B_{a_0}=A_{a_0}\), and \(D_{a_0}\) is a localization of this normal domain. The actual nonzero \(a_0\) therefore supplies the generic normal open required by the N-1 criterion for \(D\).

If each of those \(B_{\mathfrak m'}\) is N-1, that criterion makes \(D\) N-1. Finite domain descent then makes the unchanged \(S_{\mathfrak m}\) N-1. This proves why MC-\AlgebraIdentifierBreak{}STK-\AlgebraIdentifierBreak{}ERR-\AlgebraIdentifierBreak{}0721 needs only maximal ideals over the fixed \(\mathfrak m\). The source's stronger antecedent involving every maximal ideal would be sufficient, but its following proof establishes exactly the restricted set needed here. The accompanying copyedit changes ``the polynomial \ldots{} split'' to ``the polynomial \ldots{} splits''.

For one such \(\mathfrak m'\), the full product \[
 f(x)=\prod_{i=1}^{d}(x-\rho_i)\in\mathfrak m'
\] puts some factor \(u=x-\rho_i\) in \(\mathfrak m'\). Keep its exact expression \[
 u=(b-a_0\rho_i)/a_0,\qquad x=u+\rho_i,
 \qquad A[u]=A[x]=B.
\] The polynomial comparison sends \(X\) to \(U+\rho_i\) and has inverse \(U\mapsto X-\rho_i\). It preserves the original ring and local ideal; it does not replace the original fraction or lose the root factor. If \(u=0\), then \(x=\rho_i\in A\) and \(B=A\), so \(B_{\mathfrak m'}\) is normal and N-1. This is the zero case made explicit by MC-\AlgebraIdentifierBreak{}STK-\AlgebraIdentifierBreak{}ERR-\AlgebraIdentifierBreak{}0722 in the source's reused variable. For the remaining case \(u\ne0\), we prove analytic unramifiedness of \(B_{\mathfrak m'}\) by the earlier criterion applied to this specific \(u\).

The original comma-list subscripts \(\sum_{i=1,\ldots,d}\) and \(\prod_{i=1,\ldots,d}\) specify the same full index set \(\{1,\ldots,d\}\) as lower and upper limits. They are valid mathematical notation. OCC-01014 and OCC-01312 do not identify a defect, and no notation-normalizing source edit is made. In particular all repeated roots and every coefficient remain in the displayed sum and product.

\subsubsection{The linear kernel, fractional ideal and local DVR map}\label{algebra-proof-067-the-linear-kernel-fractional-ideal-and-local-dvr-map}

Retain the actual evaluation \(A[U]\to B\), \(U\mapsto u\), with kernel \(I\). For every \(a\in A\) with \(au=b'\in A\), the linear form \(aU-b'\) belongs to \(I\). These forms generate \(I\). To prove this, take its original nonzero polynomial \[
 g(U)=c_dU^d+\cdots+c_1U+c_0,\qquad c_d\ne0.
\] A constant polynomial in the kernel is zero since \(A\subset B\), so \(d\geq1\). Multiplying the actual equation \(g(u)=0\) by \(c_d^{d-1}\) yields \[
 (c_du)^d+\sum_{j=0}^{d-1}c_jc_d^{d-1-j}(c_du)^j=0.
\] This is the full monic equation in algebra.\AlgebraIdentifierBreak{}tex:\AlgebraIdentifierBreak{}30480--\AlgebraIdentifierBreak{}30496; every leading-coefficient power is retained. Thus \(c_du\) is integral over the normal \(A\) and lies in its fraction field, so \(b'=c_du\in A\). The exact subtraction is \[
 g(U)=(c_dU-b')U^{d-1}+g_1(U),
 \quad g_1\in I,\quad\deg(g_1)<d.
\] In particular the new coefficient of \(U^{d-1}\) is the old coefficient plus \(b'\); lower coefficients remain. Induction proves generation. MC-\AlgebraIdentifierBreak{}STK-\AlgebraIdentifierBreak{}ERR-\AlgebraIdentifierBreak{}0723 restores the missing plus sign in the source's displayed polynomial, as independently reported by OCC-01313 and OCC-12434.

Put \(J=A\cap uA\), an ideal of \(A\). The coefficient-at-zero map sends each generating linear form to \(-b'\), and all elements \(b'\in J\) occur because they have the form \(b'=au\) with \(a\in A\). Thus the exact mutually inverse quotient maps give \[
 B/uB=A[U]/(I,U)\cong A/J,
 \qquad [a]\longmapsto[a],\quad [h(U)]\longmapsto[h(0)].
\] The sign \(-b'\) generates the same ideal through multiplication by the explicit unit \(-1\). The Serre note, section ``Height-one intersections with all original factors retained'', applies to the same nonzero fraction \(u=(b-a_0\rho_i)/a_0\). It proves that every associated prime of \(A/J\) has height one and that there are no embedded primes, retaining all original numerator, denominator and valuation factors. Under the displayed quotient isomorphism, associated primes correspond exactly: the annihilator of an element is the inverse image of its annihilator over the quotient ring. Thus \(B/uB\), and then its localization at \(\mathfrak m'\), have no embedded primes. This supplies criterion (2).

Let \(\mathfrak q\subseteq\mathfrak m'\) be one of the associated primes of \(B/uB\), and put \(\mathfrak p=\mathfrak q\cap A\). The correspondence makes \(\mathfrak p\) an associated prime of \(A/J\), hence of height one. The ring \(A_{\mathfrak p}\) is a DVR. The localized fractional-ideal equality is \[
 J_{\mathfrak p}=A_{\mathfrak p}\cap uA_{\mathfrak p}.
\] It follows by localizing the kernel sequence defining the intersection inside the original fraction field. It is a proper ideal, since \(\mathfrak p\) is in the support of \(A/J\). If \(v_{\mathfrak p}(u)\leq0\), the second fractional ideal contains \(A_{\mathfrak p}\), contradicting properness. Hence \[
 e=v_{\mathfrak p}(u)>0,\qquad u=w\pi^e,
 \quad w\in A_{\mathfrak p}^{\times},
\] for an actual uniformizer \(\pi\), with the actual unit \(w=u\pi^{-e}\) retained. In particular \(u\in A_{\mathfrak p}\). Localizing the actual ring \(B=A[u]\) at \(A\setminus\mathfrak p\) therefore gives exactly \(A_{\mathfrak p}\). The prime \(\mathfrak q\) becomes its maximal ideal, so further localization gives \[
 B_{\mathfrak q}=A_{\mathfrak p}.
\] This proves criterion (3)(a), including the precise map justifying the source's equality. It does not infer equality merely from both rings having dimension one.

The same quotient isomorphism gives \(B/\mathfrak q\cong A/\mathfrak p\). After localizing at \(\mathfrak m'\), its exact denominator images give a local ring of this quotient of \(A\). Finite Nagata ascent preserves the quotient, and localization preserves Nagata; the result is a local Nagata domain. The local-domain theorem above makes it analytically unramified, establishing criterion (3)(b). The remaining criterion (1) is the already separated case \(u\ne0\). Hence \(B_{\mathfrak m'}\) is analytically unramified and N-1. The finite semilocal descent in the preceding section returns this result to the original \(S_{\mathfrak m}\) and then to \(S\). The original finite-type stability proposition follows.

For 46224--46245, fields and complete Noetherian local rings have already been proved Nagata. A characteristic-zero Dedekind domain is Noetherian and normal; its nonzero prime quotients are fields, and its zero-prime quotient is N-2 by the original characteristic-zero trace argument. This includes \(\mathbf Z\). Finite-type stability now gives all the displayed finite-type extensions. These implications complete the source's earlier forward reference in the infinite-polynomial example; they do not change that example's original fields or polynomial generators.

\subsubsection{The original non-Japanese DVR and the completion root test}\label{algebra-proof-067-the-original-non-japanese-dvr-and-the-completion-root-test}

At 46247--46262 retain exactly the earlier example \[
 A=\left\{\sum_{i\geq0}a_i x^i\in k[[x]]:
 [k^p(a_0,a_1,\ldots):k^p]<\infty\right\},
 \qquad [k:k^p]=\infty.
\] The complete proof that \(A\) is a DVR with uniformizer \(x\), residue field \(k\), completion \(k[[x]]\) and \(\operatorname{Frac}(A)\cap k[[x]]=A\) is retained in the Krull--Akizuki note, section ``The two original examples and their exact completions''. Keep the source's arbitrary \(f=\sum_{i\geq0}a_i x^i\notin A\) and \(R=A[f]\). Since \(f^p\in k^p[[x]]\subset A\), the original \(R\) is finite over \(A\). For every \(n\geq1\), the actual elements are \[
 g_n=\sum_{0\leq i<n}a_i x^i\in A,
 \qquad h_n=\frac{f-g_n}{x^n}=\sum_{j\geq0}a_{n+j}x^j,
 \qquad h_n^p=\sum_{j\geq0}a_{n+j}^{p}x^{pj}\in k^p[[x]]\subset A.
\] The finite set of coefficients of \(g_n\) generates a finite extension of \(k^p\), so its claimed membership follows from the definition. Each \(h_n\) is integral over \(A\), hence belongs to the closure \(R'\) in the original fraction field of \(R\). If \(R'\) were finite over \(R\), it would be finite over \(A\). The exact identities \[
 f=g_n+x^n h_n
\] put the class of \(f\) in every \(x^n(R'/A)\). The quotient \(R'/A\) is then a finite module over the local Noetherian \(A\), and the intersection of these powers is zero. Hence \(f\in A\), a contradiction. This proves the stated failure of N-1 for \(R\). If \(A\) were N-2, finite domain ascent would make \(R\) N-2 and thus N-1, also a contradiction. A DVR's only prime quotients are itself and its residue field, so the displayed equivalence of Nagata and N-2 is exact.

For 46264--46286 let \((A,\mathfrak m)\) be the original local Nagata domain of positive characteristic \(p\), \(K=\operatorname{Frac}(A)\), \(a\in A\), and \(\alpha\in\widehat A\) with \(\alpha^p=a\). The preceding local-domain theorem has already proved that \(\widehat A\) is reduced; we use it with the actual completion map. Suppose no root belongs to \(A\). If \(\beta=b/c\in K\), \(b,c\in A\), \(c\ne0\), satisfied \(\beta^p=a\), retain the full identity in \(\widehat A\) \[
 (c\alpha-b)^p=c^p\alpha^p-b^p=c^pa-b^p=0.
\] Reducedness gives \(c\alpha=b\). Faithful flatness contracts the original principal ideal exactly: \(c\widehat A\cap A=cA\), since \(A/cA\to\widehat A/c\widehat A\) is injective. Thus \(b=cd\) for some \(d\in A\), and \(\beta=d\) is the forbidden root. This avoids treating \(c\) as a unit of \(\widehat A\); equivalently the fraction comparison takes place in its total quotient ring, into which the original nonzero \(c\) remains a nonzerodivisor by flatness.

Therefore \(T^p-a\) is irreducible over \(K\), and \[
 B=A[T]/(T^p-a)\hookrightarrow K[T]/(T^p-a)
\] is a domain. The injection follows by the unique representatives of degree less than \(p\) in the monic quotient, and the target is a field. It is a finite free \(A\)-module on \(1,T,\ldots,T^{p-1}\) and is Nagata by finite ascent. Its special fibre is \[
 B/\mathfrak mB=\kappa(\mathfrak m)[T]/(T^p-\overline a).
\] The polynomial either is irreducible or is \((T-\gamma)^p\) for its unique root \(\gamma\) in the residue field, so this ring has exactly one prime. Every maximal ideal of the finite integral algebra \(B\) lies over \(\mathfrak m\), proving that \(B\) is local with maximal ideal \(\mathfrak n\). In both cases \[
 \mathfrak n^p\subseteq\mathfrak mB\subseteq\mathfrak n,
 \qquad\mathfrak n^{pk}\subseteq\mathfrak m^kB\subseteq\mathfrak n^k.
\] Thus the identity gives inverse comparisons of the maximal-adic and \(\mathfrak mB\)-adic completion systems. These are the locality and topology facts made explicit by MC-\AlgebraIdentifierBreak{}STK-\AlgebraIdentifierBreak{}ERR-\AlgebraIdentifierBreak{}0724.

Completion in the displayed finite free basis gives the actual isomorphism \[
 \widehat B=\widehat A\otimes_A B
 =\widehat A[T]/(T^p-a).
\] In it, retain \(\eta=[T]-\alpha\). The full Frobenius equation is \(\eta^p=[T]^p-\alpha^p=a-a=0\). Yet \(\eta\ne0\), because its coefficient of \([T]\) in the free basis of this monic degree-\(p\) quotient is 1, and \(p\geq2\). More explicitly, the inverse presentation maps \(T\mapsto\alpha+\eta\) and \(\eta\mapsto[T]-\alpha\) give \(\widehat B\cong\widehat A[\eta]/(\eta^p)\). The original generator \([T]\) remains \(\alpha+\eta\). Thus \(\widehat B\) is not reduced, contradicting the proved analytic unramifiedness of the local Nagata domain \(B\). The original \(a\) must have its \(p\)th root in \(A\).

\subsubsection{Report dispositions and completed receiving checks}\label{algebra-proof-067-report-dispositions-and-completed-receiving-checks}

All ten existing units MC-\AlgebraIdentifierBreak{}STK-\AlgebraIdentifierBreak{}ERR-\AlgebraIdentifierBreak{}0715--\AlgebraIdentifierBreak{}0724 are retained, with both operations of 0717. Their roles are the nonzero quotient charts, article, maximal-prime boundary before changing generators, legitimate normality reduction, missing subject, polynomial indeterminate, restricted maximal ideals, zero-generator case, missing plus sign and locality/completion comparison. The two reports about comma-list operator indices are rejected as nondefects. The ambient-field repetition in OCC-01009 is recorded as an optional clarification not integrated. The remaining received copyedits add ``by'', the enumeration colon, remove the doubled conditional and parenthesize the two coefficient rings.

The full-source reading also propagates the earlier nonzero-localization correction to the first covering assertion, restricts the transcendental case to positive transcendence degree, and repairs the singular verb ``splits''. These are source-bound corrections, with no invented received report. The generic normal opens, faithful-flat total-quotient comparison, shared localization sets, all field-generator multipliers and inverse comparisons, the exact fractional ideal and its valuation factors, every completion filtration, and the nilpotent \([T]-\alpha\) have been verified above. The Cohen, Japanese, Serre and Krull--Akizuki arguments are used only through their proved original hypotheses. No new optional extension or novelty claim is made. The next interval starts at 46293, Ascending properties; only its opening through 46322 has been read here.
